\documentclass[11pt,a4paper]{article}
\usepackage[margin=1in]{geometry}
\usepackage[T1]{fontenc}
\usepackage{lmodern}
\usepackage{xcolor}
\usepackage{microtype}
\usepackage{parskip}
\pagestyle{empty}

\definecolor{accent}{HTML}{173F5F}
\newcommand{\key}[1]{\textbf{\textcolor{accent}{#1}}}

\begin{document}

\begin{center}
{\Large\bfseries Intelligent Emergency Response \& Resource Optimization Platform}\\[0.6em]
{\large\bfseries Abstract}
\end{center}

Disasters that trigger mass-casualty events often destroy the very \key{power} and \key{communication infrastructure} needed to coordinate response. Instead of treating this as an edge case, we propose a \key{hierarchical, offline-first, event-driven architecture} that continues operating in \key{connected, degraded, and isolated modes}.

Each regional \key{Emergency Operation Center (EOC)} maintains a lightweight local broker, incident and resource state, decision engine, and human-dispatch capability, enabling immediate operation without dependence on a central service. Under normal connectivity, regional events synchronize through a durable event backbone for \key{explainable prioritization}, nationwide coordination, and continuous re-optimization. Critical incidents receive \key{instant greedy assignment}, while lower-priority incidents use \key{short-window batch optimization} to improve response time and fleet utilization.

Resources are ranked using \key{ETA, capability, capacity, hospital readiness, route safety, and regional coverage}. Fast geospatial caches shortlist candidates, while \key{transactional, version-checked reservations} prevent double dispatch. A \key{saga-based workflow} reserves, notifies, confirms, or automatically compensates when a unit fails to acknowledge.

Communication degrades gracefully from \key{4G/broadband} to \key{SMS/USSD}, \key{mesh radio}, and finally \key{store-and-forward synchronization}. Region-owned assets remain locally dispatchable, while shared scarce resources require signed ownership leases. On reconnection, urgent incidents and unresolved claims synchronize first, followed by telemetry and historical logs.

Regional isolation, circuit breakers, replicated storage, signed communication, audit trails, observability, and human override make the platform \key{scalable, secure, fault-tolerant, and explainable} even when the disaster disrupts the infrastructure it relies upon.

\end{document}
